\section*{Solution}
\begin{enumerate}
  \item L'oxydant diiode $\mathrm{I}_{2}$ oxyde le réducteur $\mathrm{S}_{2} \mathrm{O}_{3}^{2-}$ :
\end{enumerate}

$$
\begin{aligned}
\mathrm{I}_{2}+2 e^{-} & =2 \mathrm{I}^{-} \\
2 \mathrm{~S}_{2} \mathrm{O}_{3}^{2-} & =\mathrm{S}_{4} \mathrm{O}_{6}^{2-}+2 e^{-} \\
\frac{\mathrm{I}_{2}(a q)+2 \mathrm{~S}_{2} \mathrm{O}_{3}^{2-}(a q)}{} & \rightarrow 2 \mathrm{I}^{-}(a q)+\mathrm{S}_{4} \mathrm{O}_{6}^{2-}(a q)
\end{aligned}
$$

\begin{enumerate}
  \setcounter{enumi}{1}
  \item a. La quantité de diiode dans $20,0 \mathrm{~mL}$ de solution de concentration c s'écrit:
\end{enumerate}

$$
n_{1_{2}}=\left[I_{2}\right] \times V=c \times 20 \times 10^{-3}=2,0 \times 10^{-2} c \mathrm{~mol} .
$$

b. Quand le volume $v$ de solution de thiosulfate a été versé, la quantité d'ions thiosulfate introduits est :

$$
n_{\mathrm{S}_{2} \mathrm{O}_{3}^{2-}}=\left[\mathrm{S}_{2} \mathrm{O}_{3}^{2-}\right] \times V=0,10 \times v \times 10^{-3}=1,0 \times 10^{-4} v(v \text { en } \mathrm{mL}) .
$$

\begin{center}
\begin{tabular}{|l|l|l|l|l|l|}
\hline
\multicolumn{2}{|l|}{Équation} & \multicolumn{4}{|c|}{$\mathrm{I}_{2}(a q)+2 \mathrm{~S}_{2} \mathrm{O}_{3}^{2-}(a q) \rightarrow 2 \mathrm{I}^{-}(a q)+\mathrm{S}_{4} \mathrm{O}_{6}^{2-}(a q)$} \\
\hline
\multirow{2}{*}{État du système} & \multirow{2}{*}{Avancement} & \multicolumn{4}{|c|}{Quantités de matière} \\
\hline
 &  & $\mathrm{I}_{2}(a q)$ & $\mathrm{S}_{2} \mathrm{O}_{3}^{2-}(a q)$ & $\mathrm{I}^{-}(a q)$ & $\mathrm{S}_{4} \mathrm{O}_{6}^{2-}(a q)$ \\
\hline
État initial (mol) & $x=0$ & $2,0 \times 10^{2} c$ & $1,0 \times 10^{-4} v$ & 0 & 0 \\
\hline
Au cours du dosage & x & $2,0 \times 10^{-2} c-x$ & $1,0 \times 10^{-4} v-2 x$ & $2 x$ & $x$ \\
\hline
\end{tabular}
\end{center}

c. On est à l'équivalence lorsque les deux réactifs $\mathrm{I}_{2}(a q)$ et $\mathrm{S}_{2} \mathrm{O}_{3}^{2-}(a q)$ ont totalement disparu dans le bécher où s'effectue le dosage.\\
Écrivons alors l'annulation des deux quantités de matière :

$$
\left.\begin{array}{l}
n_{\mathrm{I}_{2}(u q)}=2,0 \times 10^{-2} c-x=0 \\
n_{\mathrm{S}_{2} \mathrm{O}_{6}^{2-}(a q)}=1,0 \times 10^{-4} v_{\dot{e q}}-2 x=0
\end{array}\right\} \Rightarrow 1,0 \times 10^{-4} v_{\dot{e q}}-4,0 \times 10^{-2} c=0
$$

d. La concentration du diiode est :

$$
c=2,5 \times 10^{-3} \times 12,6=3,2 \times 10^{-2} \mathbf{~ m o l} \cdot \mathrm{~L}^{-1} .
$$

\section*{Commentaires}
Écrire d'abord les demi-équations correspondant à chaque couple d'oxydoréduction.

L'avancement $x$ est ici la quantité de matière d'ions tétrathionate formés au cours de la transformation.

La définition de l'équivalence au cours d'un dosage fait partie des connaissances exigibles.



